\documentclass[]{article}
\usepackage[spanish.mexico]{babel}
\usepackage[T1]{fontenc}
\usepackage[utf8]{inputenc}
%\usepackageB.\\{lmodern}
\usepackage[a4paper]{geometry}

%\usepackage{natbib}
\usepackage{cite}


%Grafico de barras
\usepackage{pgfplots}

%Graficos e imagenes
\usepackage{graphicx}
\usepackage{subcaption}

%circuitos
\usepackage{tikz}
\usepackage[american voltages, american currents,siunitx]{circuitikz}

%Plotting

\usepackage{pgfplots}
\pgfplotsset{width=10cm,compat=1.9} 


%\title{Proyecto de Optimización de Energía}
%\author{Pablo Vivar Colina}
%\date{Mayo 2018}

%agregar portada

\begin{document}
	
	%\usepackage[top=2cm,bottom=2cm,left=1cm,right=1cm]{geometry}


\begin{titlepage}
     \begin{center}
	\includegraphics[width=0.09\textwidth]{UNAM}\Large Universidad Nacional Autónoma de México
        	\includegraphics[width=0.09\textwidth]{FI}\\[1cm]
        \Large Facultad de Ingeniería\\[1cm]
       % \Large División de Ciencias Básicas\\[1cm]
         \Large Laboratorio de Fundamentos de Control(6655)\\[1cm]
         %la clave antes era:4314
         \footnotesize Profesor: Salcedo Ubilla María Leonor Ing.\\[1cm]
        \footnotesize Semestre 2019-1\\[1cm]
        
       

        \Large Práctica No. 1\\[1cm]
        
           

\Large Introdcción MATLAB
        
         %Texto a la derecha
          \begin{flushright}
\footnotesize  Grupo 2\\[0.5cm]
\footnotesize Brigada: 4\\[0.5cm]
\footnotesize Rodrigo Adrián Martínez López\\[0.5cm]
\footnotesize Vivar Colina Pablo\\[0.5cm]
 \end{flushright}
    %Texto a la izquierda
          \begin{flushleft}
        \footnotesize Ciudad Universitaria Agosto de 2018.\\
          \end{flushleft}
         
          
        %\vfill
        %\today
   \end{center}
\end{titlepage}
 
	

	%\maketitle
	
	%\tableofcontents  % Write out the Table of Contents
	
	%\listoffigures  % Write out the List of Figures
		
	\section{Introducción}
	
	\subsection{Representación fasorial}
		
	La corriente alterna se puede representar con una flecha girando a velocidad angular $\omega$. Este elemento recibe el nombre de fasor y se representa como un número complejo.\\
	
	Su longitud coincide con el valor máximo de la tensión o corriente (según sea la magnitud que se esté representando). También se utiliza el valor RMS en lugar del valor máximo (ver transformación a fasores). En ese caso habría que dividir el valor máximo por raíz de 2.\\
	
	El ángulo (corrimiento de la señal sobre el eje horizontal) representa la fase. La velocidad de giro $\omega$ está relacionada con la frecuencia de la señal.\cite{FisicaPractica}\\

		
	\section{Objetivos}
	
    Para verificar el diagrama, se aplica al lado de alta tensión un sistema trifásico de voltajes, tomando lecturas con un voltímetro, interconectando a la vez una terminal de alta tensión y una de baja tensión.
     
	\section{Resultados}
	
	
	\begin{table}[h!]
		\centering
		\begin{tabular}{|c|c|c|}
			\hline
			Terminales Alta	& Terminales Baja  & Defasamiento  \\ \hline
		$\Delta$	& $\Delta$ & 0 $^o$  \\ \hline
			& $\lambda$  &  30 $^o$ \\ \hline
		$\lambda$ 	& $\lambda$  & 0 $^o$ \\ \hline
		$\lambda$ 	& $\Delta$ &  30 $^o$\\ \hline
		
		\end{tabular}
		\caption{Conexiones en Transformador}
	\end{table}
	
	\begin{table}[h!]
		\centering
		\begin{tabular}{|c|c|}
			\hline
		$T_2[V]$	& Medición $T_1[V]$  \\ \hline
		286.6	& $H_3-X_2=$ 56 \\ \hline
		286.8	& $H_3-X_3=$ 55.2 \\ \hline
		292.7	& $H_1-H_3=$ 109.7 \\ \hline
		280.3	& $H_2-X_2=$ 56.2 \\ \hline
		290.3	& $H_2-H_3=$ 151.6  \\ \hline
		\end{tabular}
		\caption{Relacion de conexiones}
	\end{table}
	
	Las relaciones mostradas a continuación deben cumplirse para las sitaciones donde existan $30^o$ de defasamiento.\\
	
	\begin{eqnarray}
	H_3-X_2=H_3-X_3\\
	H_3-X_2<H_1-H_3\\
	H_2-X_2<H_2-X_3\\
	H_2-X_2<H_1-H_3
	\end{eqnarray}
	
	Para el primer transformador se obtuvo lo siguiente:\\
	
	\begin{eqnarray}
	56[V]\approx 55.2[V]\\
	56[V]<109.7[V]\\
	56.2<151.6[V]\\
	56.2<109.7[V]
	\end{eqnarray}
	
	Para el segundo transformador se obtuvo lo siguiente:\\
	
	\begin{eqnarray}
	286.6[V] \approx 286.8[V]\\
	286.6[V]<292[V]\\
	280.3[V]<290.3[V]\\
	280.3<292.7[V]
	\end{eqnarray}
	
		Para el tercer transformador se obtuvo lo siguiente:\\
		
		\begin{eqnarray}
		285.2[V] \approx 284.8[V]\\
		285.2[V]<284.8[V]\\
		283.2[V]<293.3[V]\\
		280.3<292.7[V]
		\end{eqnarray}
	
	%AQUI VAMOS
	
	Es importante mencionar que las cargas lineales no deforman la forma de la onda, sin embargo las cargas no lineales sí lo hacen.\\
	
	Ejemplos de cargas lineales pueden ser una resistencia o una inductancia, y ejemplos de cargas no lineales pueden ser la luz led, semiconductores, etc.\\
    
	cd ..
	\section{Conclusiones}
	
	El objetivo de la práctica se cumplió porque logramos verificar de manera presencial el desfasamiento fasorial que se producen entre las conexiones delta y estrela que pueden presentar los transformadores.

	\section{Referencias}
	

	\bibliographystyle{plain}
    \bibliography{Referencias}
    


	
	
\end{document}

